% Part: normal-modal-logic
% Chapter: syntax-and-semantics
% Section: schemas

\documentclass[../../../include/open-logic-section]{subfiles}

\begin{document}

\olfileid{nml}{syn}{sch}

\olsection{સૂત્ર-ઢાંચા અને માન્યતા}

\begin{defn}
  \emph{સૂત્ર-ઢાંચો} એ !!{formula}sનો એવો ગણ છે જેમાં
  કોઈ મોડલ !!{formula}~$!C$ના બધા અને માત્ર બધા
  પ્રતિસ્થાપન દૃષ્ટાંતો હોય; એટલે કે,
  \[
  \Setabs{!B}{\lexists[!D_1], \dots, \lexists[!D_n] \left(!B =
  \SSubst{!C}{\subst{!D_1}{p_1}, \dots, \subst{!D_n}{p_n}} \right) }.
  \]
  !!{formula}~$!C$ને આ સૂત્ર-ઢાંચાનું \emph{લક્ષણરૂપ}
  !!{formula} કહે છે; !!{propositional
  variable}sનાં નામ બદલવા સિવાય તે અનન્ય છે.
  !!^a{formula}~$!A$ આ ગણનું સભ્ય હોય તો તેને
  સૂત્ર-ઢાંચાનો \emph{દૃષ્ટાંત} કહે છે.
\end{defn}

અણુ-ઘટકોને બદલે `$!A$', `$!B$',~\dots{} મૂકી
$!C$માંથી મળતી અધિભાષાકીય અભિવ્યક્તિ વડે
સૂત્ર-ઢાંચો દર્શાવવો અનુકૂળ છે. તેથી, દાખલા તરીકે,
`$!A$', `$!A \lif \Box!A$', `$!A \lif (!B \lif !A)$'
નીચેના સૂત્ર-ઢાંચા દર્શાવે છે. તેમનાં લક્ષણરૂપ
!!{formula}s અનુક્રમે $p$, $p \lif \Box p$, $p \lif (q
\lif p)$ છે. `$!A$' સૂત્ર-ઢાંચો \emph{બધાં}
!!{formula}sનો ગણ દર્શાવે છે.

\begin{defn}
  સૂત્ર-ઢાંચો તેના બધા દૃષ્ટાંતો નિદર્શનમાં સત્ય હોય
  ત્યારે અને માત્ર ત્યારે જ તે નિદર્શનમાં \emph{સત્ય} છે;
  અને તે દરેક નિદર્શનમાં સત્ય હોય ત્યારે અને માત્ર
  ત્યારે જ \emph{માન્ય} છે.
\end{defn}

\begin{prop}\ollabel{prop:Kvalid}
  નીચેનો સૂત્ર-ઢાંચો \Ax{K} માન્ય છે:
  \begin{equation*}
  \Box(!A \lif !B) \lif (\Box !A \lif \Box !B). \tag{\Ax{K}}
  \end{equation*}
\end{prop}

\begin{proof}
  સૂત્ર-ઢાંચાના બધા દૃષ્ટાંતો દરેક નિદર્શનના દરેક
  જગતમાં સત્ય છે તે બતાવવું છે. તેથી કોઈપણ
  $\mModel{M} = \tuple{W,R,V}$ અને $w \in
  W$ લો. શરતી વિધાન જગતમાં સત્ય છે તે બતાવવા,
  તેનો પૂર્વપક્ષ સત્ય માની તેનો ઉત્તરપક્ષ પણ સત્ય
  છે તે બતાવીએ. અહીં $\mSat{M}{\Box(!A \lif !B)}[w]$
  અને $\mSat{M}{\Box !A}[w]$ માનો. હવે
  $\mSat{M}{\Box !B}[w]$ બતાવવું છે. તો $Rww'$
  ધરાવતું કોઈપણ $w'$ લો. પ્રથમ ધારણા મુજબ
  $\mSat{M}{!A \lif !B}[w']$ અને બીજી મુજબ
  $\mSat{M}{!A}[w']$. તેથી $\mSat{M}{!B}[w']$.
  $w'$ ગમે તે હતું, તેથી $\mSat{M}{\Box !B}[w]$.
\end{proof}

\begin{prop}\ollabel{prop:Dual-valid}
  નીચેનો સૂત્ર-ઢાંચો \Dual{} માન્ય છે:
  \begin{equation*}
  \Diamond !A \liff \lnot\Box\lnot !A. \tag{\Dual}
  \end{equation*}
\end{prop}

\begin{proof}
  અભ્યાસપ્રશ્ન.
\end{proof}

\begin{prob}
  \olref[nml][syn][sch]{prop:Dual-valid} સાબિત કરો.
\end{prob}

\begin{prop}\ollabel{prop:soundMP}
  કોઈ નિદર્શનના એક જગતમાં $!A$ અને $!A \lif !B$
  સત્ય હોય તો $!B$ પણ ત્યાં સત્ય છે. તેથી માન્ય
  !!{formula}s મોડસ પોનેન્સ હેઠળ બંધ છે.
\end{prop}

\begin{prop}\ollabel{prop:valid-instances}
  !!^a{formula}~$!A$ ત્યારે અને માત્ર ત્યારે જ માન્ય છે
  જ્યારે તેના બધા પ્રતિસ્થાપન દૃષ્ટાંતો માન્ય હોય.
  બીજા શબ્દોમાં, સૂત્ર-ઢાંચો ત્યારે અને માત્ર ત્યારે જ
  માન્ય છે જ્યારે તેનું લક્ષણરૂપ !!{formula} માન્ય હોય.
\end{prop}

\begin{proof}
  ``જો'' દિશા સ્પષ્ટ છે, કારણ કે $!A$ પોતાનો જ
  પ્રતિસ્થાપન દૃષ્ટાંત છે.

  ``માત્ર જો'' દિશા સાબિત કરવા નીચેનું બતાવીએ. માનો કે
  $\mModel{M} = \tuple{W, R, V}$ મોડલ નિદર્શન છે,
  અને $!B \ident
  \SSubst{!A}{\subst{!D_1}{p_1},\dots,\subst{!D_n}{p_n}}$
  એ~$!A$નો પ્રતિસ્થાપન દૃષ્ટાંત છે. $V'(p_i) =
  \Setabs{w}{\mSat{M}{!D_i}[w]}$ મૂકી $\mModel{M'} = \tuple{W, R,
    V'}$ વ્યાખ્યાયિત કરો. ત્યારે કોઈપણ~$w \in W$ માટે
  $\mSat{M}{!B}[w]$ જો અને માત્ર જો $\mSat{M'}{!A}[w]$.
  (આ દાવાનું પ્રમાણ અભ્યાસપ્રશ્ન તરીકે છોડીએ છીએ.)
  હવે માનો કે $!A$ માન્ય છે, પણ તેનો કોઈ પ્રતિસ્થાપન
  દૃષ્ટાંત $!B$ એ~$!A$નો હોવા છતાં માન્ય નથી. તો કોઈ $\mModel{M} =
  \tuple{W, R, V}$ અને કોઈ $w \in W$ માટે
  $\mSat/{M}{!B}[w]$. પરંતુ દાવા મુજબ ત્યારે
  $\mSat/{M'}{!A}[w]$, એટલે $!A$ માન્ય નથી;
  આ વિરોધાભાસ છે.
\end{proof}

\begin{prob}
  \olref[nml][syn][sch]{prop:valid-instances}ના પ્રમાણની
  ``માત્ર જો'' દિશામાંનો દાવો સાબિત કરો. (સૂચન:
  $!A$ પર આગમન વાપરો.)
\end{prob}

જોકે, સૂત્ર-ઢાંચો કોઈ નિદર્શનમાં સત્ય હોય ત્યારે
અને માત્ર ત્યારે જ તેનું લક્ષણરૂપ સૂત્ર ત્યાં સત્ય
હોય એવું સાચું નથી. ``માત્ર જો'' દિશા અલબત્ત સાચી છે:
જો $!A$નો દરેક દૃષ્ટાંત~$\mModel{M}$માં સત્ય હોય,
તો $!A$~પોતે પણ~$\mModel{M}$માં સત્ય છે. પરંતુ
$!A$ એ~$\mModel{M}$માં સત્ય હોય છતાં $!A$નો કોઈ દૃષ્ટાંત
$\mModel{M}$ના કોઈ જગતમાં અસત્ય હોઈ શકે છે. સરળ
વિરોધી દૃષ્ટાંત માટે~$p$ લો અને માત્ર એક જગત~$w$ અને
$V(p) = \{w\}$ ધરાવતું નિદર્શન લો; તેમાં $p$
એ~$w$માં સત્ય છે. પરંતુ $\lfalse$ એ~$p$નો
દૃષ્ટાંત છે અને~$w$માં સત્ય નથી.

\begin{prob}
  નીચેના !!{formula}sમાંથી એકેય માન્ય નથી તે બતાવો:
  \begin{enumerate}
    \item[\Ax{D}:] \quad $\Box p \lif \Diamond p$;
    \item[\Ax{T}:] \quad $\Box p \lif p$;
    \item[\Ax{B}:] \quad $p \lif \Box\Diamond p$;
    \item[\Ax{4}:] \quad $\Box p \lif \Box \Box p$;
    \item[\Ax{5}:] \quad $\Diamond p \lif \Box \Diamond
      p$.
  \end{enumerate}
\end{prob}

\begin{table}[t]
    \centering
    \begin{tabular}{| l || l |}
      \hline
      {\emph{માન્ય સૂત્ર-ઢાંચા}} & {\emph{અમાન્ય સૂત્ર-ઢાંચા}} \\
      \hline\hline
      $\Box(!A \lif !B) \lif (\Diamond !A \lif \Diamond !B)$
      & $\Box (!A \lor !B) \lif (\Box !A \lor \Box !B)$ \\
      $\Diamond (!A \lif !B) \lif (\Box !A \lif \Diamond
      !B)$
      & $(\Diamond !A \land \Diamond !B) \lif \Diamond (!A
      \land !B)$\\
      $\Box (!A \land !B) \liff (\Box !A \land \Box !B)$
      & $!A \lif \Box !A$ \\
      $\Box !A \lif \Box (!B \lif !A)$
      & $\Box \Diamond !A \lif !B$ \\
      $\lnot \Diamond !A \lif \Box (!A \lif !B)$
      & $\Box \Box !A \lif \Box !A$ \\
      $\Diamond (!A \lor !B) \liff (\Diamond !A \lor
      \Diamond !B)$
      & $\Box \Diamond !A \lif \Diamond \Box !A$. \\
      \hline
    \end{tabular}
    \caption{માન્ય અને અમાન્ય સૂત્ર-ઢાંચા.}
    \ollabel{tab:valid-invalidSchemas}
\end{table}
\sourcecorrection{OLMD-010}{સ્થિર અંગ્રેજી મૂળના કોષ્ટક-શીર્ષકમાં
``and (or?) invalid'' એવો અનિર્ણીત સંપાદકીય કૌંસ છે.
બંને સ્તંભની માન્યતા અને અમાન્યતા સાન્ત વિરોધી
નિદર્શનો સામે તપાસીને ગુજરાતી શીર્ષકમાં તે કૌંસ દૂર કર્યો છે.}

\begin{prob}%\ollabel{ex:in/validSchemas}
  \olref[nml][syn][sch]{tab:valid-invalidSchemas}ના
  પ્રથમ સ્તંભના સૂત્ર-ઢાંચા માન્ય છે અને બીજા સ્તંભના
  સૂત્ર-ઢાંચા માન્ય નથી તે સાબિત કરો.
\end{prob}

\begin{prob}
  નીચેના સૂત્ર-ઢાંચા માન્ય છે કે અમાન્ય, તે નક્કી કરો:
  \begin{enumerate}
  \item $(\Diamond !A \lif \Box !B) \lif (\Box !A \lif \Box
    !B)$;
  \item $\Diamond(!A \lif !B) \lor \Box(!B \lif !A)$.
  \end{enumerate}
\end{prob}

\begin{prob}
  નીચેના દરેક સૂત્ર-ઢાંચા માટે એવું નિદર્શન $\mModel{M}$
  શોધો કે $\mModel{M}$માં તે !!{formula}નો દરેક દૃષ્ટાંત સત્ય હોય:
  \begin{enumerate}
  \item $p \lif \Diamond\Diamond p$;
  \item $\Diamond p \lif \Box p$.
  \end{enumerate}
\end{prob}

\end{document}
